\section*{Chapter \ref{ch:m2:settlement-risk} --- Settlement Risk}

\begin{solution}{exo:m2:settlement-risk:1}
Selling yen for dollars: from the yen deadline at 22:00 the evening before to the
dollar confirmation at 21:00, 23 hours. Selling dollars for yen: with these hours
the yen arrive before the dollar payment becomes irrevocable, so there is no window;
in practice earlier instruction deadlines can create one.
\end{solution}

\begin{solution}{exo:m2:settlement-risk:2}
Gross: the whole principal of the currency it paid, since the counterparty's leg
never comes. Under PvP: only the cost of replacing the trade at the new rate, since
its own payment is not made either.
\end{solution}

\begin{solution}{exo:m2:settlement-risk:3}
The first bank pays EUR~100 million net, and the dollars are netted in the same way:
one payment in each currency instead of twenty.
\end{solution}

\begin{solution}{exo:m2:settlement-risk:4}
A: the bank sells EUR for USD~100 million net. B: sells yen for USD~100 million. C:
sells euros for sterling, 100 million, and dollars for sterling, 100 million. D:
sells the emerging currency for USD~30 million. USD~430 million in all.
\end{solution}

\begin{solution}{exo:m2:settlement-risk:5}
Each leg is paid only when both are, so members must have each currency they owe
at CLS in time, without first receiving what they are owed in the other; that is a
funding need across currencies and hours. Settling many trades together lets
payments in and out offset, so only net pay-ins are needed.
\end{solution}

\begin{solution}{exo:m2:settlement-risk:6}
Many emerging-market currencies are not settled in PvP systems, their payment
systems' hours may not overlap with the dollar's, and smaller institutions often
lack direct or affordable access to PvP or netting services: their trades settle
gross, over correspondent accounts.
\end{solution}

\begin{solution}{exo:m2:settlement-risk:7}
Gross: USD~330, 930, 1\,330 and 730 million at 05:00, 08:00, 14:00 and 17:00;
netted: 130, 330, 430 and 230 million.
\end{solution}

\begin{solution}{exo:m2:settlement-risk:8}
The market value of the forwards is the replacement cost, the exposure before
settlement. On the settlement date the exposure is the whole principal of the
currency paid until the other arrives, often hundreds of times larger, and it
recurs every value date. Settlement exposures need their own limits.
\end{solution}

\begin{solution}{pb:m2:settlement-risk:1}
\textbf{1.} USD~1.53 billion of purchases; a peak of USD~1.33 billion.
\textbf{2.} From 13:00, when the dollar payments become irrevocable, to 16:00, when
the euros and sterling are confirmed: sales of dollars are then exposed while the
earlier sales of euros and yen still await their dollars.
\textbf{3.} The sale of yen for USD~250 million to B: 23 hours.
\textbf{4.} USD~330 million: B's yen sale and D's sale of the emerging currency.
\textbf{5.} USD~700 million: the euros sold for dollars (400 million, awaiting the
dollars) and the dollars sold for euros (300 million, awaiting the euros).
\textbf{6.} As in exercise 4: USD~100, 100, 200 and 30 million with A, B, C and D.
\textbf{7.} USD~430 million.
\textbf{8.} USD~100 million.
\textbf{9.} A netting agreement enforceable in the counterparty's insolvency, so
that a liquidator cannot claim the gross amounts owed to the failed bank while not
paying those it owes.
\textbf{10.} USD~30 million: the net sale of the emerging currency, because that
currency is not settled by PvP.
\textbf{11.} USD~30 million.
\textbf{12.} Only replacement risk, a fraction of the principal, since their trades
settle payment against payment.
\textbf{13.} Fees and membership or access costs, the liquidity to pay in each
currency on time, and the operations to meet the deadlines.
\textbf{14.} By a PvP arrangement for that currency if one exists, by asking D to pay
first, by collateral or a limit, or by settling through a bank that offers the
currency with loss protection.
\textbf{15.} Because the counterparty has no access to PvP, the currency is not
supported, or the cost is judged too high for the trade.
\textbf{16.} As a credit limit on the peak settlement exposure per value date,
sized to the counterparty's credit quality and the bank's capital, and checked
before trades are accepted.
\textbf{17.} It showed that FX trades carry principal risk, not only price risk, and
led to settlement limits, netting and in time PvP.
\textbf{18.} That about USD~1.4 trillion a day still depends on each payer trusting
that the other currency will arrive: a large, recurring exposure concentrated in
certain currencies and counterparties.
\textbf{19.} Named result: \emph{the unsettled day}: the book's peak exposure is
USD~1.33 billion settled gross, USD~430 million netted, and USD~30 million with PvP
for every currency but one.
\textbf{20.} It is the risk of paying one currency and not receiving the other, and
\omterm{def:m2:settlement-risk:pvp}{payment versus payment} removes it.
\end{solution}

\begin{solution}{iq:m2:settlement-risk:1}
The risk that one leg of an FX trade is paid irrevocably, in its own time zone,
while the other, due later elsewhere, is not received because the counterparty
fails, as in Herstatt's closure in 1974. CLS settles both legs across its books at
the same moment, or neither, so no principal is exposed.
\iqlookfor{the time-zone mechanism and PvP.}
\end{solution}

\begin{solution}{iq:m2:settlement-risk:2}
\omterm{def:m2:settlement-risk:risk}{Settlement risk} is the loss of the full principal when one leg is paid and the
other never arrives. Replacement risk is the loss from having to redo an unsettled
trade at a worse rate after the counterparty fails: a fraction of principal,
related to price moves.
\iqlookfor{principal versus market move.}
\end{solution}

\begin{solution}{iq:m2:settlement-risk:3}
For each trade due, find the time its sold currency's payment becomes irrevocable
and the time its bought currency is confirmed received; the exposure to the
counterparty at time $t$ is the sum of bought amounts with $t$ in that window,
after netting where agreements apply and excluding PvP-settled trades; extend by
unconfirmed and failed receipts; compute it by value date and in real time as
confirmations arrive.
\iqlookfor{windows from the payment schedules, netting, and confirmations.}
\end{solution}

\begin{solution}{iq:m2:settlement-risk:4}
PvP covers only some currencies and time windows; some counterparties lack access
directly or through a member; joining or using it costs fees and liquidity; some
trades (same-day, late, certain products) do not fit its schedule; and some banks
judge the risk acceptable within limits.
\iqlookfor{access, coverage, cost and timing.}
\end{solution}

\begin{solution}{iq:m2:settlement-risk:5}
By as much as the two parties' flows offset: for dealers trading both ways all day,
by a large share; for one-way clients, little. It needs trades on the same value
date and currencies, a legally enforceable netting agreement, and operations that
agree the net amounts in time.
\iqlookfor{offsetting flows and the legal basis.}
\end{solution}

\begin{solution}{iq:m2:settlement-risk:6}
Maintain, per counterparty and value date, the projected exposure profile from
booked trades, the payment schedules and netting; on each new trade, compute the
profile with it and compare its peak with the limit before acceptance, in the
pre-trade path; update as confirmations arrive; escalate or reject when a limit
would be breached; test against historical days and failure scenarios.
\iqlookfor{pre-trade check on the projected peak, updated by confirmations.}
\end{solution}
